\documentclass[11pt]{amsart}

\usepackage[T1]{fontenc}
\usepackage{lmodern}
\usepackage{microtype}
\usepackage{mathtools,amssymb,amsthm}
\usepackage[hidelinks]{hyperref}

\hypersetup{
  pdftitle={The first r+1 columns of the cubic-dimension spectral sequence of a
            hypercube Vietoris-Rips complex}
}

\newtheorem{theorem}{Theorem}[section]
\newtheorem{lemma}[theorem]{Lemma}
\newtheorem{proposition}[theorem]{Proposition}
\newtheorem{corollary}[theorem]{Corollary}
\theoremstyle{definition}
\newtheorem{problem}[theorem]{Problem}
\theoremstyle{remark}
\newtheorem{remark}[theorem]{Remark}

\DeclareMathOperator{\diam}{diam}
\DeclareMathOperator{\cdim}{cdim}
\DeclareMathOperator{\Ind}{Ind}
\DeclareMathOperator{\sgn}{sgn}
\newcommand{\kk}{\Bbbk}
\newcommand{\Hy}{\mathcal H}
\newcommand{\Fc}{\mathcal F}
\newcommand{\bl}{\bullet}
\numberwithin{equation}{section}

\title[The first $r+1$ columns of a hypercube Rips spectral sequence]
{The First $r+1$ Columns of the Cubic-Dimension Spectral Sequence
 of a Hypercube Vietoris--Rips Complex}

\date{}

\subjclass[2020]{55N31, 05E18, 55T05, 05C10}
\keywords{Vietoris--Rips complex, hypercube graph, cubic dimension,
spectral sequence, hyperoctahedral group}

\begin{document}

\begin{abstract}
Let $X^{n,r}$ be the Vietoris--Rips complex of the hypercube graph $Q_n$ at
scale $r$, filtered by cubic dimension, and let $E^\bl_{\bl,\bl}$ be the
resulting spectral sequence over a field of characteristic zero, all as in
Galetto, Monta\~no and Wellner (arXiv:2606.20784). Problem~8.2 of that paper
asks for a proof that $E^1_{p,0}\cong\Fc_p$, the $p$-th term of the cellular
chain complex of the $n$-cube, and that $E^1_{p,q}=0$ for $q\neq0$, for all
$0\le p\le r$. We prove both, for every $n>r\ge0$. Because $p\le r$ the Rips
condition is vacuous on the vertex set of a $p$-cell, so the $p$-th column is,
exactly and not merely up to homotopy, the relative chain complex of a simplex
modulo the union of its $2p$ half-cube faces; an affine marginal map identifies
that pair $\Hy_p$-equivariantly with $(I^p,\partial I^p)$. For $r\le3$ both
assertions are already proved in the source paper.
\end{abstract}

\maketitle

\section{The problem}\label{sec:problem}

Let $Q_n$ be the hypercube graph: vertex set $V(Q_n)=\{0,1\}^n$, with the
Hamming distance $d$. For an integer $r\ge0$ the \emph{Vietoris--Rips complex}
$X^{n,r}$ is the simplicial complex whose faces are the non-empty subsets
$S\subseteq V(Q_n)$ with $\diam(S)\le r$. The \emph{cubic dimension} of such an
$S$ is
\[
  \cdim(S)=\#\{\,i\in[n] : \text{$x_i\neq y_i$ for some $x,y\in S$}\,\},
\]
the number of coordinates in which $S$ is not constant, and $X^{n,r}_p$ denotes
the subcomplex of faces with $\cdim\le p$. This is an increasing filtration of
$X^{n,r}$ by subcomplexes. Writing $C^{n,r}_{d,p}$ for the span, over a fixed
field $\kk$, of the faces of dimension $d$ and cubic dimension \emph{exactly}
$p$, the associated homology spectral sequence has
\[
  E^0_{p,q}=C^{n,r}_{p+q,\,p},
  \qquad
  E^1_{p,q}=H_{p+q}\bigl(C^{n,r}_{p+\bl,\,p}\bigr).
\]
The hyperoctahedral group $\Hy_n=(\mathbb Z/2)^n\rtimes S_n$ acts on $Q_n$ by
isometries, hence on everything above; all isomorphisms below are isomorphisms
of $\Hy_n$-representations. Finally $\Fc_\bl$ is the cellular chain complex of
the $n$-cube $[0,1]^n$, so that $\Fc_p$ has as basis the $p$-cells of the cube
and $\dim\Fc_p=\binom{n}{p}2^{n-p}$.

All of the above is the setting of Galetto, Monta\~no and Wellner \cite{GMW};
for other homological invariants of $X^{n,r}$ see \cite{BEG}. Section~8 of
\cite{GMW}, ``Open problems'', poses the following.

\begin{problem}[\cite{GMW}, Problem 8.2]\label{prob:82}
Let $n>r\ge0$. Show that for all integers $p$ with $0\le p\le r$ we have
\begin{enumerate}
\item[(i)] $E^1_{p,0}\cong\Fc_p$, where $\Fc_\bl$ is the cellular chain complex
  of the $n$-cube; and
\item[(ii)] $E^1_{p,q}=0$ for all $q\neq0$.
\end{enumerate}
\end{problem}

The statement as posed quantifies explicitly only over $p$; the ambient
quantifiers above are the ones we read off \cite{GMW}, whose main theorem is
stated for $n>r$ and which fixes $\kk$ to be of characteristic zero. The labels
$\{\lambda;\mu\}$ for irreducible $\Hy_m$-representations are those of
\cite{GMW}; $\{0;1^m\}$ is the determinant character of $\Hy_m$, i.e.\ the sign
of the underlying permutation times the product of the signs.

\section{The result}\label{sec:result}

\begin{theorem}\label{thm:main}
Let $n>r\ge0$ and let $\kk$ be a field of characteristic zero, as in
\cite{GMW}. Then for every $p$ with $0\le p\le r$:
\begin{enumerate}
\item[(i)] $E^1_{p,0}\cong\Fc_p$ as $\Hy_n$-representations; and
\item[(ii)] $E^1_{p,q}=0$ for all $q\neq0$.
\end{enumerate}
\end{theorem}

The characteristic hypothesis is inherited from the two results quoted from
\cite{GMW} in Propositions~\ref{prop:free} and~\ref{prop:F}, and from the
labelling $\{0;1^p\}$; the local input we supply, Lemma~\ref{lem:pair}, is
characteristic-free, so over an arbitrary field the same argument still gives
$E^1_{p,q}=0$ for $q\neq0$ and $\dim E^1_{p,0}=\dim\Fc_p$, with $\Hy_p$ acting
by the determinant on each local summand.

For $r\le3$ both assertions of Problem~\ref{prob:82} are already proved in
\cite{GMW} itself. We make no claim of priority beyond that, and none about
literature we have not read. The proof below needs no computer;
Section~\ref{sec:scope} describes the program that accompanies this note.

\subsection*{Two results quoted from \cite{GMW}, not reproved here}

\begin{proposition}[\cite{GMW}]\label{prop:free}
For all $p$ and $q$,
\[
  E^1_{p,q}\;\cong\;
  \Ind_{S_{n-p}\times\Hy_p}^{\Hy_n}
  \bigl(\,\{n-p\}\otimes H_{p+q}(C^{p,r}_{p+\bl,\,p})\,\bigr),
\]
where $\{n-p\}$ is the trivial representation of $S_{n-p}$.
\end{proposition}

\begin{proposition}[\cite{GMW}]\label{prop:F}
For all $i$,
$\Fc_i\cong\Ind_{S_{n-i}\times\Hy_i}^{\Hy_n}(\{n-i\}\otimes\{0;1^i\})$, and
$\{0;1^i\}$ is one-dimensional. In particular
$\dim\Fc_i=\binom{n}{i}2^{n-i}$.
\end{proposition}

These two are exactly the reduction we import. Induction is exact, so
Proposition~\ref{prop:free} says that the $p$-th column is determined by the
case $n=p$ alone. Combining the two propositions:

\begin{corollary}\label{cor:reduce}
Fix $r\ge0$ and $0\le p\le r$. If
\[
  H_{k}\bigl(C^{p,r}_{p+\bl,\,p}\bigr)=
  \begin{cases}
    \{0;1^p\} & k=p,\\
    0 & k\neq p
  \end{cases}
  \qquad\text{as $\Hy_p$-representations,}
\]
then assertions \emph{(i)} and \emph{(ii)} of Problem~\ref{prob:82} hold for
that $p$ and \emph{every} $n>r$.
\end{corollary}

Only this implication is used below. It is a \emph{reduction}, not a proof: it
converts Problem~\ref{prob:82} into a statement about one local complex, and it
supplies neither that complex's homology nor its character. Those are
Sections~\ref{sec:object} and~\ref{sec:termwise}.

\section{The collapse}\label{sec:object}

Fix $p$ with $0\le p\le r$ and put $n=p$; write $V=V(Q_p)=\{0,1\}^p$, so
$|V|=2^p$.

\begin{lemma}\label{lem:collapse}
Since $\diam(Q_p)=p\le r$, the Rips condition is vacuous on $V$: every
non-empty subset of $V$ is a face of $X^{p,r}$. Hence
\[
  \Delta:=X^{p,r}=\text{the full simplex on the $2^p$ vertices $V$},
\]
and a face $S\subseteq V$ has $\cdim(S)=p$ if and only if $S$ is
\emph{full-support}, i.e.\ $S$ is non-constant in every coordinate. Therefore
\[
  U:=X^{p,r}_{p-1}
   =\bigcup_{i\in[p]}\;\bigcup_{\varepsilon\in\{0,1\}}\Delta(F_{i,\varepsilon}),
  \qquad F_{i,\varepsilon}:=\{x\in V: x_i=\varepsilon\},
\]
the union of the $2p$ \emph{half-cube simplices} $\Delta(F_{i,\varepsilon})$
\textup{(}each a full simplex on $2^{p-1}$ vertices\textup{)}, and there is an
equality of chain complexes of $\Hy_p$-representations
\begin{equation}\label{eq:collapse}
  C^{p,r}_{p+\bl,\,p}\;=\;C_{p+\bl}(\Delta,U).
\end{equation}
\end{lemma}

\begin{proof}
Any two vertices of $Q_p$ are at Hamming distance $\le p\le r$, so every
non-empty $S\subseteq V$ satisfies $\diam(S)\le r$: this is the whole of the
first assertion, and it is where the hypothesis $p\le r$ is used. Next,
$\cdim(S)<p$ means exactly that $S$ is constant in some coordinate $i$, i.e.\
$S\subseteq F_{i,\varepsilon}$ for some $\varepsilon$, which is the stated
description of $U$; and $\cdim(S)=p$ means $S$ is full-support. So the faces of
$\Delta$ of cubic dimension exactly $p$ are precisely the faces of $\Delta$ not
in $U$, and $C^{p,r}_{d,p}=C_d(\Delta)/C_d(U)=C_d(\Delta,U)$ with the
simplicial differential on both sides. Every step is $\Hy_p$-equivariant
because $\Hy_p$ preserves the Hamming distance and hence cubic dimension.
\end{proof}

Equation \eqref{eq:collapse} is an \emph{equality}, not a homotopy
equivalence, and it is the whole idea: the hypothesis $p\le r$ deletes all
metric content, leaving an object with $2p$ pieces. At $p=4$, for instance,
$\Delta$ has $2^{16}-1=65535$ faces and $U$ has $1760$, so
$C^{4,4}_{4+\bl,4}$ has $63775$ generators; the next section obtains its
homology from $U$, with no rank computation.

\section{The local pair}\label{sec:termwise}

\begin{lemma}[the marginal map]\label{lem:pair}
There is an $\Hy_p$-equivariant homotopy equivalence of pairs
\[
  \bigl(|\Delta|,|U|\bigr)\;\simeq\;\bigl(I^p,\partial I^p\bigr),
  \qquad I=[0,1],
\]
where $\Hy_p$ acts on $I^p$ as the affine symmetry group of the cube. Hence for
all $k$
\[
  H_k(\Delta,U)\cong H_k(I^p,\partial I^p)=
  \begin{cases}\kk & k=p,\\ 0 & k\neq p,\end{cases}
\]
and $\Hy_p$ acts on $H_p(\Delta,U)\cong\kk$ by the determinant of the linear
part, i.e.\ by the character $\{0;1^p\}$.
\end{lemma}

\begin{proof}
Realise $|\Delta|$ as the set of probability measures $\mu$ on $V=\{0,1\}^p$,
the face of $|\Delta|$ containing $\mu$ in its interior being
$\operatorname{supp}\mu$. Let
\[
  f:|\Delta|\longrightarrow I^p,\qquad
  f(\mu)=\bigl(\mathbb E_\mu x_1,\dots,\mathbb E_\mu x_p\bigr),
\]
an affine map, and let $s:I^p\to|\Delta|$ be the Bernoulli section
$s(t)=\bigotimes_{i=1}^p\mathrm{Bern}(t_i)$, so that $f\circ s=\mathrm{id}$.

First, $f^{-1}(\partial I^p)=|U|$ \emph{exactly}: $f(\mu)_i=\mu(\{x_i=1\})$
lies in $\{0,1\}$ if and only if $\operatorname{supp}\mu\subseteq F_{i,0}$ or
$\operatorname{supp}\mu\subseteq F_{i,1}$, i.e.\ if and only if
$\operatorname{supp}\mu$ fails to be full-support in coordinate $i$; and by
Lemma~\ref{lem:collapse} that is exactly membership of $U$.

Second, the straight-line homotopy
$H(\mu,u)=(1-u)\mu+u\,s(f(\mu))$ takes values in $|\Delta|$ (a convex
combination of probability measures) and satisfies, because $f$ is affine and
$f\circ s=\mathrm{id}$,
\[
  f\bigl(H(\mu,u)\bigr)=(1-u)f(\mu)+u\,f(\mu)=f(\mu)
  \qquad\text{for all }u .
\]
So $f$ is constant along $H$; in particular $H$ preserves
$f^{-1}(\partial I^p)=|U|$ at every time. Thus $H$ is a deformation retraction
of the pair $(|\Delta|,|U|)$ onto $(s(I^p),s(\partial I^p))$, and $s$ is a
homeomorphism onto its image with inverse $f$.

Third, equivariance. For $g=(s,\sigma)\in\Hy_p$ acting on $V$ by
$(gx)_{\sigma(i)}=x_i+s_{\sigma(i)}$, one has
$\mathbb E_{g\mu}x_{\sigma(i)}=\mathbb E_\mu x_i$ if $s_{\sigma(i)}=0$ and
$=1-\mathbb E_\mu x_i$ if $s_{\sigma(i)}=1$. Hence $f$ intertwines the
$\Hy_p$-action on $|\Delta|$ with the affine action
$t\mapsto\bigl(t_{\sigma^{-1}(i)}\ \text{or}\ 1-t_{\sigma^{-1}(i)}\bigr)$ of
$\Hy_p$ on $I^p$, which is the full affine symmetry group of the cube; $s$ is
equivariant for the same two actions, and $H$ is then equivariant as well.

Finally $H_k(I^p,\partial I^p)=\widetilde H_{k-1}(S^{p-1})$ is $\kk$ for $k=p$
and $0$ otherwise (for $p=0$ read $\partial I^0=\emptyset$ and
$H_0(\mathrm{pt},\emptyset)=\kk$), and an affine symmetry $g$ of $I^p$ acts on
the fundamental class $[I^p,\partial I^p]$ by $\det$ of its linear part, which
for $g=(s,\sigma)$ is $\sgn(\sigma)(-1)^{s_1+\dots+s_p}$. That is the
determinant character of $\Hy_p$, i.e.\ $\{0;1^p\}$ in the labelling of
\cite{GMW}.
\end{proof}

\begin{proof}[Proof of Theorem~\ref{thm:main}]
By \eqref{eq:collapse} and Lemma~\ref{lem:pair},
\[
  H_k\bigl(C^{p,r}_{p+\bl,\,p}\bigr)=H_k(\Delta,U)=
  \begin{cases}\{0;1^p\}&k=p,\\ 0&k\neq p,\end{cases}
\]
for every $0\le p\le r$. Corollary~\ref{cor:reduce} converts this into (i)
and (ii) for every $n>r$.
\end{proof}

\section{Scope, and the accompanying program}\label{sec:scope}

Three limits. First, the labelling $\{\lambda;\mu\}$ of $\Hy_m$-irreducibles is
used as in \cite{GMW}, over a field of characteristic zero;
Lemma~\ref{lem:pair} itself, and hence Theorem~\ref{thm:main} read as
``$\dim$, and the determinant character'', is characteristic-free. Second, the
truncation $p\le r$ is essential to the argument: at $p=r+1$ the minimal
full-support subsets of $V(Q_p)$ are the antipodal pairs, at distance $p>r$, so
the Rips condition is no longer vacuous and Lemma~\ref{lem:collapse} does not
apply. Third, Theorem~\ref{thm:main} is a statement about the terms
$E^1_{p,q}$; we prove nothing here about the differential $d^1$, nor about the
later pages.

Theorem~\ref{thm:main} is proved above without a computer. A program
\texttt{verify.py}, with its recorded output \texttt{verify.output.txt},
accompanies this note; the recorded run reports $144$ checks, all passing. It
rebuilds the local complexes $C^{p,r}_{p+\bl,\,p}$ from the definitions of
Section~\ref{sec:problem}, checks \eqref{eq:collapse} as an identity of basis
sets for $p\le4$, and confirms the conclusion of Lemma~\ref{lem:pair} --- one
class, in degree $p$, on which $\Hy_p$ acts by the determinant --- for $p\le4$,
computing the homology of the $63775$-generator complex $C^{4,4}_{4+\bl,4}$
both through $U$ and directly by exact elimination on its fifteen boundary
matrices. It also reproduces the dimensions \cite{GMW} publishes for its own
scale-$2$ and scale-$3$ computations. It is Python~3.9+, standard library only;
all arithmetic is exact integer arithmetic, and ranks are taken modulo a prime,
which bounds the characteristic-zero Betti numbers from above, the
field-independent alternating sum then pinning them.

The program reproves neither Proposition~\ref{prop:free} nor
Proposition~\ref{prop:F}, does not model the homotopy of Lemma~\ref{lem:pair},
and enumerates no cell with $p=r\ge5$, which would require $2^{32}$ subsets;
its output lists each check it runs and each thing it does not cover, and some
of those checks concern objects the argument above does not use. The proof in
Sections~\ref{sec:object} and~\ref{sec:termwise} is uniform in $n$, $r$ and $p$
and does not depend on any computation.

\begin{thebibliography}{9}

\bibitem{GMW}
F.~Galetto, J.~Monta\~no, and Z.~Wellner,
\emph{Homology of Vietoris--Rips complexes of hypercube graphs via group
actions},
arXiv:2606.20784v1, 2026.
\href{https://arxiv.org/abs/2606.20784}{arXiv:2606.20784}.

\bibitem{BEG}
M.~Bendersky, N.~Elia, and J.~Grbi\'c,
\emph{Cohomological properties of the Vietoris--Rips complex of a hypercube
graph},
arXiv:2605.00705v2, 2026.
\href{https://arxiv.org/abs/2605.00705}{arXiv:2605.00705}.

\end{thebibliography}

\end{document}
